\documentclass{article}

\usepackage{amsfonts}
\usepackage{mathtools}
\usepackage{amsthm}
\usepackage{mathrsfs}
\usepackage{amsmath}
\usepackage{graphicx}
\usepackage{hyperref}
\usepackage{stackengine}
\usepackage[mathscr]{eucal}
\usepackage{listings}
\usepackage{algorithm}
\usepackage{algpseudocode}
\usepackage{float}
\usepackage{xcolor}
\usepackage{url}
\newtheorem{theorem}{Theorem}
\theoremstyle{definition}\newtheorem{definition}{DEF}
\theoremstyle{definition}\newtheorem{example}{Ex}
\theoremstyle{lemma}\newtheorem{lemma}{Lemma}
\theoremstyle{corrolary}\newtheorem{corrolary}{Cor}
\setlength{\parindent}{0pt}
\definecolor{codebg}{RGB}{248,248,248}
\definecolor{codekeyword}{RGB}{0,76,153}
\definecolor{codecomment}{RGB}{0,128,96}
\definecolor{codestring}{RGB}{163,21,21}
\definecolor{codenumber}{RGB}{102,102,102}
\lstdefinestyle{pythoncode}{
    language=Python,
    basicstyle=\ttfamily\small,
    keywordstyle=\color{codekeyword}\bfseries,
    commentstyle=\color{codecomment}\itshape,
    stringstyle=\color{codestring},
    numberstyle=\scriptsize\color{codenumber},
    numbers=left,
    numbersep=10pt,
    stepnumber=1,
    showstringspaces=false,
    breaklines=true,
    keepspaces=true,
    columns=fullflexible,
    tabsize=4,
    frame=single,
    framerule=0pt,
    rulecolor=\color{codebg},
    backgroundcolor=\color{codebg}
}
\lstset{style=pythoncode}
\DeclareMathOperator*{\argmax}{arg\,max}
\DeclareMathOperator*{\argmin}{arg\,min}

\begin{document}

Likelihood is one of the big ways that we understand statistics.
As we have seen, there are a variety of estimators that can have properties such as unbiasedness and consistency.
However, we may wish to know where these estimators come from.
Likelihood is an unifying theory that allows us to to createa  new estimators for different distributions and examine their properties.

\begin{definition}
Let $X$ be a random variable with parameter $\theta$ and pdf $f$.
Say the likelihood of $\theta$ is
\begin{equation}
{\mathrm L}(\theta | x) = f(x | \theta )
\end{equation}
\end{definition}

\begin{lemma}
Let $X_1, X_2, ..., X_n$ be n iid random variables with parameter $\theta$ and likelihood ${\mathrm L}$.
Then the joint likelihood is
\begin{equation}
{\mathrm L}(\theta | x_1, x_2, ..., x_n) = \prod_{k=1}^n {\mathrm L}(\theta | x_k)
\end{equation}
\end{lemma}
\begin{proof}
First, notice the random variables are iid, so the value of $\theta$ is the same across all $1 \leq k \leq n$.
Also, all variables have the same PDF, $f$.
Thus,
\begin{align*}
{\mathrm L}(\theta | x_1, x_2, ..., x_n)
&= f(x_1, x_2, ..., x_n | \theta ) \\
&= \prod_{k=1}^n f(x_k | \theta ) \\
&= \prod_{k=1}^n {\mathrm L}(\theta | x_k)
\end{align*}
\end{proof}

Notice that likelihood is not a probability distribution.
Despite the fact that it is equal to the pdf, we use the term likelihood to refer to when the observations are fixed and the parameter is variable.
INdeed, notice that if $X\sim Exp(\lambda )$, then
\begin{align*}
\int_\mathbb{R} {\mathrm L} d\lambda
&= \int_0^\infty \lambda e^{-\lambda x}d\lambda \\
&= \frac{-\lambda}{x} e^{-\lambda x} -\frac{1}{x^2} e^{-\lambda x}|_0^\infty \\
&= \frac{1}{x^2} \\
&\neq 1
\end{align*}

\begin{definition}
Let $X$ be a random variable with a parameter $\theta$ and likelihood ${\mathrm L}$.
Say the log likelihood is
\begin{equation}
\ell (\theta | x) = \ln {\mathrm L}(\theta | x)
\end{equation}
\end{definition}

\begin{lemma}
Let $X_1, X_2, ..., X_n$ be $n$ iid random variables with parameter $\theta$ and log likelihood $\ell$.
Thenthe joint likelihood is
\begin{equation}
\ell (\theta | x_1, x_2, ..., x_n) = \sum_{k=1}^n \ell (\theta | x_k)
\end{equation}
\end{lemma}
\begin{proof}
From the definition of log likelihood, we have
\begin{align*}
\ell(\theta | x_1, x_2, ..., x_n )
&= \ln {\mathrm L} (\theta | x_1, x_2, ..., x_n) \\
&= \ln \prod_{k=1}^n {\mathrm L} (\theta | x_n) \\
&= \sum_{k=1}^n \ln {\mathrm L} (\theta | x_k) \\
&= \sum_{k=1}^n \ell (\theta | x_k)
\end{align*}
\end{proof}

\begin{definition}\label{mle}
Let $X$ be a random variable with parameter $\theta$ from parameter space $\Theta$.
Say the maximum likelihood estimator (MLE) is the estimator
\begin{equation}
\hat{\theta} = \argmax_{\theta \in \Theta} {\mathrm L} (\theta | x)
\end{equation}
\end{definition}

\begin{lemma}
If $\hat{\theta}$ is the MLE then $\hat{\theta}$ maximizes the log likelihood.
\end{lemma}
\begin{proof}
Notice that $\ln u$ is a one-to-one increasing function, so if $\hat{\theta}$ maximizes ${\mathrm L} (\theta | x)$ it also maximizes $\ell (\theta | x) = \ln {\mathrm L} (\theta | x)$.
\end{proof}

\begin{example}
Let $X\sim Poisson(\lambda)$ and $\{x_k\}$ be an SRS of size $n$.
So the $x_k$ are independent and the likelihood is $\ell(\theta|x)=x\ln \lambda - \lambda - \ln x!$.
Thus, the joing likelihood is
\begin{align*}
\ell (\theta | \{x_k\})
&= \sum_{k=1}^n\left( x_k\ln \lambda - \lambda - \ln x_k! \right) \\
&= \sum_{k=1}^n x_k \ln \lambda - n\lambda -\sum_{k=1}^n \ln x_k!
\end{align*}
Thus,
\begin{equation*}
\frac{\partial \ell}{\partial \lambda} = \frac{\sum_{k=1}^n x_k}{\lambda} - n
\end{equation*}
Thus, the MLE is
\begin{equation*}
\hat{\lambda} = \frac{n}{\sum_{k=1}^n x_k}
\end{equation*}
\end{example}

\begin{example}
Let $X\sim \mathcal{N}(\mu, \sigma^2)$ and $\{x_k\}$ be an SRS.
THen the log likelihood is
\begin{equation}
\ell (\mu, \sigma^2 | \{x_k\}) = \frac{-n}{2}\ln (2\pi) -\frac{n}{2} \ln (\sigma^2) -\frac{1}{2} \sum_{k=1}^n \frac{(x_k - \mu)^2}{\sigma^2}
\end{equation}
So
\begin{equation}
\frac{\partial \ell}{\partial \mu} = \sum_{k=1}^n \frac{-(x_k - \mu)}{\sigma^2}
\end{equation}
Thus,
\begin{equation*}
\hat{\mu} = \frac{1}{n} \sum_{k=1}^n x_k = \bar{x}
\end{equation*}
Also,
\begin{equation*}
\frac{\partial \ell}{\partial (\sigma^2)} = \frac{-n}{2\sigma^2} + \sum_{k=1}^n \frac{(x_k - \mu)^2}{2(\sigma^2)^2}
\end{equation*}
so
\begin{equation*}
\hat{\sigma^2} = \frac{1}{n}\sum_{k=1}^n (x_k - \bar{x})^2
\end{equation*}
Notice that the MLE of hte mean $\hat{\mu}$ is the sample mean $\bar{x}$ and is unbiased and consistent.
However, the MLE for the variance $\hat{\sigma^2}$ is not the sample variance $s^2$, and is biased.
\end{example}

\begin{example}
Let $X \sim Unif(0, \theta)$ and $\{x_k\}$ be an SRS.
Notice $f(x|\theta)=\frac{1}{\theta}$ for $0\leq x \leq \theta$.
So $\ell (\theta | \{x_k\})=-n\theta$ for $\theta > x_1, x_2, ..., x_n$.
Thus, the MLE is $\hat{\theta}=\max_k \{x_k\}$
\end{example}

\begin{definition}
Let $X$ be a random variable with parameter $\theta \in \Theta$  and log likelihood $\ell$.
Say the Fisher information matrix is
\begin{equation}
I_X(\theta) = \mathbb{E}\left[ \nabla_\theta \ell(\theta | x) \nabla_\theta^T \ell l(\theta | x) \right]
\end{equation}
\end{definition}

\begin{theorem}
\begin{equation}
I_X(\theta) = \mathbb{E}\left[ -\nabla_\theta^2 \ell(\theta | x)\right]
\end{equation}
\end{theorem}
\begin{proof}
We will look at a one dimensional parameter for simplicity.
THe argument can be applied on an element-by-element basis for a general parameter.
Notice that $\ell'(\theta | x) = \frac{f'(x|\theta)}{f(x|\theta)}$.
Thus,
\begin{equation*}
\ell''(\theta|x) = \frac{f(x|\theta)f''(x|\theta) - (f'(x|\theta))^2}{(f'(x|\theta))^2}
\end{equation*}
Next, notice
\begin{align*}
\mathbb{E}\left[ \frac{f''(x|\theta)f(x|\theta)}{(f(x|\theta))^2}\right]
&= \int_\mathbb{R} \frac{f''(x|\theta)}{f(x|\theta)} dF \\
&= \int_\mathbb{R} \frac{f'(x|\theta)}{f(x|\theta)} f(x|\theta) dx \\
&= \int_\mathbb{R} f''(x|\theta) dx \\
&= \frac{\partial^2}{\partial \theta^2} \left( \int_\mathbb{R} f(x|\theta)dx \right) \\
&= \frac{\partial^2}{\partial \theta^2}\left( 1 \right) \\
&= 0
\end{align*}
Thus,
\begin{align*}
\mathbb{E}\left[ -\ell''(\theta|x)\right]
&= \mathbb{E}\left[ \frac{f''(x|\theta)f(x|\theta)}{(f(x|\theta))^2}\right] - \mathbb{E}\left[\frac{(f'(\theta|x))^2}{(f(\theta|x))^2} \right]\\
&= 0 + \mathbb{E}\left[ -(\ell'(\theta|x))^2\right] \\
&= I_X(\theta)
\end{align*}
\end{proof}

\begin{theorem}
Let $X_1, X_2, ..., X_n$ be iid random variables with Fisher information matrix $I_X(\theta)$.
Then the FIM of $\sum_i X_i$ is
\begin{equation}
I(\theta)=n I_X(\theta)
\end{equation}
\end{theorem}
\begin{proof}
Let $\ell$ be the log likelihood of $x_i$.
Then, the joint log likelihood is $\ell(\theta | x_1, x_2, ..., x_n)=n\sum_i \ell(\theta | x_i)$.
Thus, the Fisher information is
\begin{align*}
I(\theta)
&= \mathbb{E} \left[ -\nabla^2 \ell (\theta | x_1, x_2, ..., x_n) \right] \\
&= \mathbb{E}\left[ \sum_i -\nabla^2 \ell (\theta | x_i )\right] \\
&= \sum_i \mathbb{E}\left[ -\nabla^2 \ell ( \theta | x_i) \right] \\
&= \sum_i I_X(\theta) \\
&= nI_X(\theta)
\end{align*}
\end{proof}

\begin{example}
Let $T\sim Exp(\lambda)$.
Then $\ell(\lambda | t) = -\lambda t + \ln \lambda$.
So $\ell''(\lambda | t) = \frac{-1}{\lambda^2}$.
Thus, $I_T(\lambda)=\frac{1}{\lambda^2}$.
\end{example}

\begin{example}
Let $X\sim \mathcal{N}(\mu, \sigma^2)$.
Then $\ell(\mu, \sigma^2 | x) = \frac{-1}{2} \ln (2\pi ) - \frac{1}{2}\ln \sigma^2 -\frac{(x-\mu)^2}{2\sigma^2}$.
So $\nabla \ell = \left( \frac{x-\mu}{\sigma^2} , \frac{(x-\mu)^2}{2(\sigma^2)^2} -\frac{1}{2\sigma^2}\right)$.
Thus, the Hessian is
\begin{equation*}
\nabla^2 \ell =
\begin{bmatrix}
\frac{-1}{\sigma^2} & \frac{x-\mu}{\sigma^2} \\
\frac{x-\mu}{\sigma^2} & \frac{-(x-\mu)^2}{(\sigma^2)^3} +\frac{1}{2(\sigma^2)^2}
\end{bmatrix}
\end{equation*}
Clearly, $\mathbb{E}\left[\frac{x-\mu}{\sigma^2}\right]=0$ and $\mathbb{E}\left[ -\left( \frac{-1}{\sigma^2}\right)\right]=\frac{1}{\sigma^2}$.
Finally, consider
\begin{align*}
\mathbb{E}\left[ \frac{(X-\mu)^2}{(\sigma^2)^3} -\frac{1}{2(\sigma^2)^2} \right]
&= \frac{1}{(\sigma^2)^2} \mathbb{E}\left[ \frac{(X-\mu)^2}{\sigma^2}\right] - \frac{1}{2(\sigma^2)^2} \\
&= \frac{1}{(\sigma^2)^2} - \frac{1}{2(\sigma^2)^2} \\
&= \frac{1}{2\sigma^4}
\end{align*}
Thus, the Fisher information is
\begin{equation*}
I_X(\mu, \sigma^2) =
\begin{bmatrix}
\frac{1}{\sigma^2} & 0 \\
0 & \frac{1}{2\sigma^4}
\end{bmatrix}
\end{equation*}
Suppose instead we view $\sigma$ as a parameter.
Then $\frac{\partial^2 \ell}{\partial \sigma^2} = \frac{1}{\sigma^2} - \frac{3(x-\mu)^2}{\sigma^4}$.
So the fisher information is $I_X(\sigma)=\frac{2}{\sigma^2}$.
\end{example}

\begin{theorem}
Let $X$ be a random variable with parameter $\theta\in \Theta$ and let $\{x_k\}$ be an SRS.
Let $\hat{\theta}$ be the unique MLE on the interior of $\Theta$ and let the log likelihood $\ell (\theta | x)$ be twice differentiable.
Then $\hat{\theta}$ is asymptotically normal with asymptotic variance $I_X(\theta)^{-1}$.
\end{theorem}
\begin{proof}
First, let $f(x|\theta)$ be the PDF.
Because $\{x_k\}$ is an SRS, $X_k$ are iid random variables.
Notice that $\frac{1}{n}\ell(\hat{\theta}|x_1, x_2, ..., x_n)=\frac{1}{n}\sum_k \ell(\hat{\theta}|x_k)$ is the maximum.
Then, $\frac{1}{n} \nabla \ell (\hat{\theta} | x_1, x_2, ..., x_n) = 0$.
Also, by the central limit theorem, $\frac{1}{n} \ell(\omega | x_1, x_2, ..., x_n) {P \atop \rightarrow} \mathbb{E}[f(x|\omega)]$.
Let $\theta^*\in \Theta$ be the true value of $\theta$, that is the value of $\theta$ such that $X\sim F(\theta^*)$.
Then
\begin{align*}
\frac{1}{n} \ell (\omega | x_1, x_2, ..., x_n) - \frac{1}{n} \ell (\theta^* | x_1, x_2, ..., x_n)
&{P \atop \rightarrow} \mathbb{E}\left[ \ln f(x|\omega) - \ln f(x|\theta^*)\right] \\
&= \mathbb{E}\left[ \ln \frac{f(x|\omega)}{f(x|\theta^*)}\right] \\
&\leq \ln \mathbb{E} \left[ \frac{f(x|\omega)}{f(x|\theta^*)}\right] \\
&= \ln \int_\mathbb{R} f(x|\omega) dx \\
&= 0
\end{align*}
Thus, $\ell(\omega | x_1, x_2, ..., x_n) {P \atop \rightarrow} \ell (\theta^* | x_1, x_2, ..., x_n)$.
Thus, for any $\omega \in \Theta$, $\frac{-1}{n} \nabla^2 \ell(\omega | x_1, x_2, ..., x_n) {P \atop \rightarrow} I_X(\theta)$. \\
Next, notice $\mathbb{E}[\nabla \ell(\theta|x)]=0$ and $Var(\nabla\ell (\theta|x)) = I_X(\theta)$.
Then, by the central limit theorem, $\frac{1}{n} \nabla \ell(\theta | x_1, x_2, ..., x_n) {D\atop \rightarrow} \mathcal{N}(0, I_X(\theta^*))$.
Also, by Taylor's Theorem, 
\begin{equation*}
\nabla(\hat{\theta} | x_1, x_2, ..., x_n) = \nabla \ell(\theta^* | x_1, x_2, ..., x_n) + (\hat{\theta} - \theta^*)^T \nabla^2 \ell (\psi | x_1, x_2, ..., x_n)
\end{equation*}
for some $\psi \in \Theta$.
Notice that $(- \nabla^2 \ell(\psi | x_1, x_2, ..., x_n))^{-1} {D\atop \rightarrow} nI_X(\theta^*)$
Thus, by Slutsky's theorem, 
\begin{align*}
\sqrt{n} (\hat{\theta} - \theta^*) 
&= \sqrt{n} (-\nabla^2 \ell(\psi | x_1, x_2, ..., x_n))^{-1} \nabla \ell (\theta^* | x_1, x_2, ..., x_n) \\
&= (\frac{-1}{n} \nabla^2 \ell(\psi | x_1, x_2, ..., x_n))^{-1} (\frac{1}{\sqrt{n}} \nabla \ell (\theta^* | x_1, x_2, ..., x_n)) \\
&{D \atop \rightarrow} \mathcal{N}(0, 1)
\end{align*}
\end{proof}

\begin{definition}\label{mvue}
Let $\theta$ be a parameter.
Say the minimum variance unbiased estimator (MVUE) is the unbiased estimator $\omega$ that has least variance of any unbiased estimator for $\theta$.
\end{definition}

\begin{theorem}
\textbf{Cramer-Rao Loewr Bound:}
Let $X$ be a random variable with parameter $\theta$ and let $\omega$ be an unbiased estimator of $\theta$.
Then
\begin{equation}
Var(\omega) \geq (nI_X(\theta))^{-1}
\end{equation}
\end{theorem}
\begin{proof}
First, let $z(\theta | x) = \frac{\partial \ell}{\\partial \theta}$.
Then $z(\theta | x_1, x_2, ..., x_n)=\sum_k z(\theta|x_k)$.
Notice that 
\begin{equation*}
Cov(\omega, z) \leq \sqrt{Var(\omega)Var(z)}
\end{equation*}
Next, because $\omega$ is unbiased,
\begin{align*}
\theta
&= \mathbb{E}[\omega] \\
&= \int_\mathbb{R} \omega dF
\end{align*}
Thus, taking the derivative gives
\begin{align*}
1 
&= \int_\mathbb{R} \omega z dF
\end{align*}
\end{proof}

\end{document}
